% !TeX root = ../papers/nc-p3.tex
\section{両向き二元点での四次接錐と残余階数15因子の全次数正則性}
\label{sec:double-binary-residual-regularity}
\begin{theorem}\label{thm:double-binary-residual-regularity}
定理\ref{thm:binary-opposite-determinant}の残余階数15因子の幾何学的一般点は
標準14641型quadratic Koszul左右AS正則である。向きを一つ移した因子にも同じ結論が成り立つ。
これらの一般窓は安定で、二十個の被約点を持ち、最小周期4である。
さらに、各向きの因子を二つの片側対称因子で切ったうち、以下で指定する側の交叉の
幾何学的一般点も全次数で正則となり、反対階数は10である。
\end{theorem}
\begin{proof}[二つの二元軌跡の横断面]
定理\ref{thm:binary-reverse-orientation}の $d=x$, $l=x-1$, $c_3=x^3+2$ の有理点を取る。
ここでは二行は独立で $\gcd(\ell,c)=d$ なので、グラフは有限で窓は安定Frobeniusである。
両向きを $B_1,B_2$ とする。各指標の隣接Jacobi核次元は有理数上で $45,16,16,0$、
両混合指標の二次障害はともに非零である。従って定理\ref{thm:binary-component-crossing}の
高さによる証明がこの点にも適用できる。滑らかな15次元枝上に $a,b$ を取り、その両対称交叉
$H=(a=b=0)$ は滑らかな13次元である。以下の次元では文字基底の軌道を除く。

各 $B_i$ の被約な芽がこの点で滑らかなことを確かめる。
向きを固定した二元点では、反対切断の一方の核は、ねじれ三次曲線 $C$ を含む二次式の鉛筆である。
その完全交叉は $C$ と残余直線 $s$ の和である。実際二つの二次式は共通因子を持たず、
三次曲線が二重に現れる次数はない。残余次数1で、完全交叉は純次元である。
この中の直線は $s$ 一つだけで、直線の入射はこのファイバーで非分岐である。
$s$ の一般点で完全交叉は滑らかなので、
$N_{s/\mathbb P^3}\to\mathcal O_s(2)^2$ は一般点で単射、従って大域切断の核が零だからである。
階数8が一定な二元軌跡上で鉛筆を束とし、その中の直線をGrassmann多様体上で取る。
入射は固有で、標本の近くで有限である。特殊ファイバーが被約な一点なので、
中山の補題で関数環は底の代数の商となり、被約な二元軌跡を覆うことからその上で同型である。
残余三次曲線も代数的に回復する。直線の二つの一次式に関して鉛筆を二行二列の一次式で書き、
その行列式を鉛筆に加えると残余曲線のイデアルになる（直線との残余完全交叉の式）。
従って滑らかなねじれ三次曲線の族を得る。曲線上の射影座標はエタール局所で選べる。
その二次像の線型包と階数1の形式の根基から三平面を戻す操作は
定理\ref{thm:binary-apolar-rank-test}の証明の定階数の線型演算である。
そこで積の二元座標と第四文字を選べば、Gram行列と八次形式も同じく代数的に戻る。
従って滑らかな二元パラメータと文字枠による表示は、これらの基底選択の自由度を除いて局所逆を持つ。
特に各被約像は滑らかな12次元である。この議論は他方の向きにも使える。

元の表示の八次形式に $(3,1,4,1,5,9,2,6,5)$ を加える一次方向を $u$ とし、
逆表示の八次形式で同じ変更をして元の基底に戻した方向を $v$ とする。
十次・八次形式の微分に二文字の軌道微分を加えた二空間の階数はそれぞれ43、
その和は44である。また第一空間に $v$、第二空間に $u$ を一本加えるといずれも44となる。
小行列は素数1009で非零であり、上界は上の次元と両対称軌道の次元31である。
よって $B_1,B_2$ は $H$ 内で互いに横断し、それぞれに含まれる選んだ二方向が
その共通部分への横断を与える。以下、$u,v$ をこの接基底の座標にも用いる。
この二方向を残した $H$ の滑らかな二次元切片では $B_1$ の接線が $v=0$、
$B_2$ の接線が $u=0$ である。混合方向 $a,b$ を加えた同変な四次元切片を取る。
\end{proof}
\begin{proof}[余階数二の残余行列]
前節と同じく交代六列と六行から $a,b$ をくくって得た行列を $M$ とする。
その四ブロックは二交換で不変である。
\path{computations/verify_double_binary_transversal.py}は $M_0=M(0)$ の階数が
有理数上で14であること、$u,v$ に沿う微分がそれぞれ
$\ker M_0\to\operatorname{coker}M_0$ 上で零になることを有理数上で検算する。
計算式は定理\ref{thm:binary-opposite-determinant}の四ブロック微分である。
$a,b$ に沿う一次微分も、二交換の符号により零である。
従って可逆な14次小行列で消去した二行二列のSchur補は四次元切片上で二次から始まり、
$D=\det M$ の消失次数は4以上である。

良い素数1009で四次斉次項は単元倍を除いて
\[
               (664v^2+546b^2)(691u^2+570a^2)
\]
となる。計算を説明する。$a,b$ 両方の速度を非零に取り、不変方向の速度を
$X=uX_1+vX_2$ とおく。隣接Schur方程式の二次を指標ごとに解き、
対称な二次補正を選んで三次の混合障害を消す。三次の混合解と四次の $--$ 解までを
用いて $M=M_0+tM_1+t^2M_2+O(t^3)$ を計算する。
行列をくくるときには $a,b$ の各速度も割り、定数行列と核基底を全方向で共通に取る。
$L,R$ を左・右核の基底列、$J_0$ を $M_0$ の一般逆として、
Schur補の二次項は $L^t(M_2-M_1J_0M_1)R$ である。
\path{computations/verify_double_binary_residual_quartic.py}はその左上成分が零、
右上・左下成分が表示の両二次式となることを検査する。
二交換で不変な二次式の単項式は $u^2,uv,v^2,a^2,b^2$ の五つだけであり、
コードはこれらを分離する方向で偏極する。符号空間の線型補正は
滑らかな枝の同変な陰関数消去で選べる。混合接の自由項は基底変更と混合座標の選択にあたり、
行列式の最低項には単元倍しか与えない。自由な対称二次補正についても
$H$ の接空間全体で左右核への微分がこの素数で零と検算しているので還元値は変わらない。
全ての消去に単元小行列を用い、標数零での三次以下の消失は上記の有理数計算で確認済みである。
従って標数零の最低次斉次項は四次で、その還元が表示式となる。
\end{proof}
\begin{proof}[両対称正則点への特殊化]
15次元枝上、$D$ の局所既約因子 $F$ の幾何学的一般点が非正則であると仮定する。
有限Frobenius開部分であるから、その因子内に全次数で正則な点が一つでもあれば、
各内部次数でのbar複体の特殊化により一般点も全次数で正則となる。
従って $F$ の $H$ への制限の零点は非正則な点だけからなる。
$D|_H$ は上の四次項から恒等零でない。有限な両対称Frobenius窓の
非正則なものは二方向の二元軌跡だけである。各 $B_i$ の一般点では
定理\ref{thm:binary-opposite-determinant}により $D|_H$ は単純に消える。
従って $F|_H$ は単元倍を除き $B_1,B_2$ の一方の局所式、または両者の積である。
この議論は完備局所環でもよい。各 $B_i$ は滑らかで、全次数の特殊化は
局所環の基底変換後にも各次数の有限複体で行える。

二交換は $H$ を点ごとに固定する。$F|_H$ に現れる $B_i$ に沿う消失が単純だから、
同じ $B_i$ を含む別の局所因子はなく、二交換は $(F)$ を保つ。
生成元は二交換で平均して不変にできる。実際 $H$ 上で元の非零な制限が変わらないため、
各単元倍率は閉点で1である。四次元切片に制限すると、
$B_i$ が一つなら $F$ の最低項は $u$ または $v$ の定数倍である。
二つなら最低項は二次で、不変かつ $a=b=0$ への制限が $uv$ の非零定数倍となる。

これらはいずれも四次項の因子にならない。標数1009の代数閉包で
表示した四次式は相異なる四つの一次式の積となる。$b$ の符号は $v,b$ の二因子を交換し、
$a$ の符号は $u,a$ の二因子を交換する。従って不変な一次因子はなく、
不変な二次因子の $a=b=0$ への制限は $u^2$ または $v^2$ の倍数である。
標数零で仮定した最低項の因数分解は、係数体上で1009を延長した付値を取って、
各因子をGauss付値零に規格化すればこの素数にも還元できる。
四次式全体の $a=b=0$ への制限がこの素数でも非零なので、どの因子の制限も
還元で零にならず、上で求めた $u,v,uv$ のいずれかがそのまま残る。これは矛盾である。
従ってこの両向き点を通る残余因子はいずれも正則な両対称窓に特殊化でき、
幾何学的一般点が全次数で正則である。各向きの一般二元点で一意な残余因子は
その二元軌跡の閉包を含むので、この点にも来る。既知の有限被約性・安定性・周期と合わせて主張を得る。

同様に $a=0$ に制限すると最低項は $u^2(664v^2+546b^2)$ の定数倍である。
$B_1$ を含む因子で非正則性を仮定すれば、不変な最低項の制限は $v$ または $uv$ の倍数を要する。
$b$ の符号を保つ因子でこれは不可能である。従って $B_1$ から来る残余因子のこの片側切片も
一般に正則である。$b=0$ では $B_2$ について同じである。
これらの一般点は両対称でなく、二元点で欠けていた中央支持3の点も閉条件により一般点では欠ける。
有限片側対称で階数9以下の正則窓の判定
（系\ref{cor:all-defective-one-sym-finite}と定理\ref{thm:finite-rank-nine-all-marks}）から
階数10を得る。残余因子上の任意の特殊値の正則性はここでは一般点への特殊化からは従わない。
\end{proof}
